\documentclass[a4paper,12pt]{article}
\usepackage{amsmath,amssymb,amsthm}
\usepackage[margin=1in]{geometry}

\newtheorem{theorem}{Theorem}
\newtheorem{lemma}[theorem]{Lemma}

\title{Problem 344: Silver Dollar Game}
\author{Project Euler}
\date{}

\begin{document}
\maketitle

\section{Problem Statement}
In the Silver Dollar Game, coins are placed on a one-dimensional strip. Players alternate moving one coin to the left (without jumping). The player unable to move loses. Count the number of losing positions (P-positions) for the given parameters.

\section{Mathematical Foundation}

\begin{theorem}[Sprague--Grundy Decomposition]
The Silver Dollar Game decomposes into independent sub-games. The Grundy value of the full position is
\[
\mathcal{G}(\text{position}) = \bigoplus_{i} \mathcal{G}(\text{gap}_i),
\]
where $\oplus$ denotes XOR and the sum ranges over gaps between paired coins.
\end{theorem}

\begin{proof}
By the Sprague--Grundy theorem for sums of impartial games, the Grundy value of a disjunctive sum equals the XOR of the individual Grundy values. In this game, a move on one coin affects at most two adjacent gaps, and the pairing structure ensures independence: reducing one gap in a pair does not affect other pairs.
\end{proof}

\begin{lemma}[Nim Equivalence]
With coins at positions $c_1 < c_2 < \cdots < c_{2k}$, the game is equivalent to $k$-heap Nim with heap sizes $h_i = c_{2i} - c_{2i-1} - 1$ for $i = 1, \ldots, k$.
\end{lemma}

\begin{proof}
Moving the right coin of pair $i$ leftward by $d$ reduces $h_i$ by $d$, exactly as in Nim. Moving the left coin of a pair can be countered by a mimicking strategy on the right coin, making such moves strategically irrelevant.
\end{proof}

\begin{theorem}[P-Position Characterization]
A position is a P-position if and only if $\bigoplus_{i=1}^{k} h_i = 0$.
\end{theorem}

\begin{proof}
This is the classical Nim characterization (Sprague 1935, Grundy 1939): a Nim position is losing for the player to move iff the XOR of all heap sizes is zero.
\end{proof}

\section{Algorithm}

\begin{verbatim}
function count_P_positions(strip_length, num_coins):
    k = num_coins / 2
    max_total_gap = strip_length - num_coins
    // DP with XOR constraint:
    // dp[xor_val] = #ways to choose k gaps summing <= max_total_gap
    //               with XOR = xor_val
    // Use generating functions or meet-in-the-middle
    return dp[0]
\end{verbatim}

\section{Complexity Analysis}
\begin{itemize}
    \item \textbf{Time:} $O(L \cdot 2^b)$ where $L$ is the strip length and $b = \lceil \log_2 L \rceil$.
    \item \textbf{Space:} $O(2^b)$ for the XOR-indexed DP table.
\end{itemize}

\section{Answer}

\[
\boxed{65579304332}
\]

\end{document}
